\begin{trapbox}

	\begin{description}[style=nextline, leftmargin=2.8em, labelsep=0.5em, itemsep=0.35em]
		\item[\textbf{\textcolor{CTrap}{易错点一（$\Delta>0$漏写）}}]
			使用两不等根条件时，常忘记判别式应大于0，误将$\Delta\ge0$作为条件，导致允许重根进入。
		\item[\textbf{\textcolor{CTrap}{正确理解（$\Delta>0$必要性）：}}]
			两个不等实根必须$\Delta>0$；若$\Delta=0$，则为重根，不满足不等要求。
		\item[\textbf{\textcolor{CTrap}{错因分析（不等根误解）：}}]
			认为“两个实根”包括相等情况，忽略了“不等”要求；或没有注意题目表述。
	\end{description}

	\medskip
	\begin{description}[style=nextline, leftmargin=2.8em, labelsep=0.5em, itemsep=0.35em]
		\item[\textbf{\textcolor{CTrap}{易错点二（对称轴漏检）}}]
			对于两不等根，只列出$\Delta>0$和$af(m)>0,af(n)>0$，忽略检查对称轴是否在区间内，导致可能两根不在区间内。
		\item[\textbf{\textcolor{CTrap}{正确理解（对称轴必在区间内）：}}]
			两根关于对称轴对称，若对称轴不在区间内，则不可能两根都在区间内。
		\item[\textbf{\textcolor{CTrap}{错因分析（对称性缺失）：}}]
			不理解对称轴与根的位置关系，或仅记忆条件不全面。
	\end{description}

	\medskip
	\begin{description}[style=nextline, leftmargin=2.8em, labelsep=0.5em, itemsep=0.35em]
		\item[\textbf{\textcolor{CTrap}{易错点三（符号分类错误）}}]
			在$af(m)>0$中，将不等式直接当作$f(m)>0$（假定$a>0$），忽视$a<0$时需反向。
		\item[\textbf{\textcolor{CTrap}{正确理解（统一符号处理）：}}]
			$af(m)>0$等价于$f(m)$与$a$同号，应根据$a$的符号具体确定$f(m)$的正负。
		\item[\textbf{\textcolor{CTrap}{错因分析（未按$a$分类）：}}]
			习惯假设开口向上（$a>0$），缺乏分类讨论习惯。
	\end{description}

\end{trapbox}
